% Positive-Z camera convention and pixel-domain Jacobian.
\clearpage
\section*{Verification case: a projection Jacobian has units and a domain}
\noindent In a positive-depth camera frame, let
\[
u=f_x\frac XZ+c_x,\qquad
v=f_y\frac YZ+c_y,\qquad Z>0.
\]
These equations map metric camera coordinates into pixel coordinates.
The local Jacobian is
\[
J=\begin{pmatrix}
f_x/Z&0&-f_xX/Z^2\\
0&f_y/Z&-f_yY/Z^2
\end{pmatrix},
\]
whose entries have units of pixels per metre if \(X,Y,Z\) are in metres
and \(f_x,f_y\) are focal lengths in pixels.
It is valid only for nonzero positive depth under the chosen
camera convention.

For \(f_x=f_y=800\), \(c_x=320\), \(c_y=240\) and
\((X,Y,Z)=(2,1,4)\), the image coordinate is \((720,440)\)
and
\[
J=\begin{pmatrix}200&0&-100\\0&200&-50\end{pmatrix}.
\]
The small camera-space displacement
\(\delta=(0.01,-0.02,0.04)\) has predicted pixel displacement
\(J\delta=(-2,-6)^T\).
The \emph{nonlinear} camera formula evaluated at the perturbed point
instead gives approximately \((-1.9802,-5.9406)\).
The discrepancy is a local linearization error, not automatically
a programming defect.

\paragraph{Verification protocol.}
For fixed positive depth, compare analytic Jacobian columns against
central finite differences as the perturbation scale decreases;
record where floating-point cancellation reverses the improvement.
Repeat nearer to the camera, where terms scale as \(1/Z\) and
\(1/Z^2\); reject zero and negative depth as outside this
derivation's domain. Keep the camera sign, pixel-axis and
column-vector conventions explicit so this mathematics can later
be translated into a renderer without accidental convention changes.
